% ============================================================
% Recreational Marijuana Legalization and Drug-Overdose
% Mortality: A Cross-Border Distance Design
% Formatted to Economic Inquiry (Wiley/WEAI) submission requirements.
% Compile: pdfLaTeX -> BibTeX -> pdfLaTeX -> pdfLaTeX (Overleaf: Recompile)
% ============================================================
\documentclass[11pt]{article}

\usepackage[margin=1in]{geometry}
\usepackage{setspace}
\usepackage{amsmath,amssymb}
\usepackage{graphicx}
\usepackage{booktabs}
\usepackage{caption}
\usepackage[authoryear,round]{natbib}
\usepackage[hidelinks]{hyperref}
\usepackage{enumitem}
\usepackage{titlesec}

\doublespacing
\titleformat{\section}{\normalfont\bfseries\large}{\thesection.}{0.5em}{}
\titleformat{\subsection}{\normalfont\bfseries}{\thesubsection}{0.5em}{}

\title{}
\author{}
\date{}

\begin{document}

% ============================================================
% TITLE PAGE
% ============================================================
\begin{titlepage}
\begin{center}
\vspace*{1.5cm}
{\LARGE \bfseries Recreational Marijuana Legalization and\\
Drug-Overdose Mortality: A Cross-Border\\
Distance Design\par}
\vspace{0.3cm}
{\large \itshape Evidence from Never-Legalizing-State Counties, 2020--2023\par}
\vspace{1.2cm}
{\large Author name(s) withheld for review\par}
\vspace{0.3cm}
{\normalsize Working Paper --- July 2026\par}
\end{center}

\vspace{1.2cm}
\noindent\textbf{Abstract} \\[4pt]
\noindent We test whether proximity to a newly legal recreational-
marijuana market reduces drug-overdose mortality, using two
co-primary treatments for counties in 32 never-legalizing states,
2020--2023: the continuous change in distance to the nearest legal
border, and a binary indicator for counties whose distance fell below
50 miles (``switchers''). Switcher counties see overdose deaths fall
3.5 more per 100,000 than controls ($p=0.030$), but the continuous-
distance specification is not significant ($p=0.103$). A supplementary
check using marijuana and alcohol arrest rates at the same threshold
reproduces a companion paper's alcohol-arrest complementarity signal
robustly across both continuous specifications. We map the legal-
market's staggered rollout across four distinct legalization waves and
interpret the overdose result as suggestive, not confirmed. \\[10pt]
\noindent\textbf{JEL Classification:} I18, I12, K42 \\
\noindent\textbf{Keywords:} recreational marijuana; opioid overdose
mortality; cross-border spillover; drug policy;
difference-in-differences

\vfill
\noindent\textbf{Abbreviations} \\[2pt]
\noindent
\begin{tabular}{@{}ll@{}}
DiD    & Difference-in-differences \\
FIPS   & Federal Information Processing Standard (county code) \\
NCHS   & National Center for Health Statistics \\
SE     & Standard error \\
SRS    & Summary Reporting System (FBI arrest data) \\
VSRR   & Vital Statistics Rapid Release \\
\end{tabular}

\vspace{0.5cm}
\noindent{\footnotesize A separate Disclosure Statement accompanies this
submission per Economic Inquiry policy and is not reproduced in this file.}
\end{titlepage}

% ============================================================
% 1. INTRODUCTION
% ============================================================
\section{Introduction}

Does living close to a legal recreational-marijuana market change
drug-overdose mortality for residents who never cross the border
themselves? The question sits inside a specific, contested empirical
literature on whether marijuana access substitutes for opioids.
\citet{Bachhuber2014} found early evidence that state medical-marijuana
laws were associated with fewer opioid overdose deaths.
\citet{Shover2019} showed the association reverses sign once later
years of data are included, and \citet{PowellPaculaJacobson2018} found
the effect, where present, is specific to dispensary-based legal
designs rather than medical-marijuana laws generally.
\citet{MathurRuhm2022} replicate and extend this line of work using
generalized difference-in-differences and event-study models on
state-level legalization timing, and conclude that legal retail
dispensaries -- medical or recreational -- are, if anything, associated
with \emph{higher} opioid mortality, casting further doubt on earlier
protective-effect estimates. Most recently, \citet{Denkyirah2026} use a
state-level two-way-fixed-effects design focused specifically on
\emph{recreational} legalization and find a robust decline of
approximately 3.5 opioid overdose deaths per 100,000 -- a magnitude we
will return to, since our own switcher-county estimate lands in almost
exactly the same place despite a completely different identification
strategy.

Every paper above identifies off the \emph{timing} of state policy
adoption. None uses \emph{geography}: whether a county's proximity to
an external legal market, in a state that has not itself legalized,
predicts its own overdose trajectory. This is the design
\citet{HaoCowan2020} developed for marijuana-possession arrests
(finding a positive spillover in border counties of non-legal
neighboring states) and that a companion project extends to
marijuana/alcohol arrest substitution and, separately, to housing
markets. To our knowledge, no paper combines this cross-border-
distance design with drug-overdose mortality as the outcome -- the
contribution of this paper.

We report two treatments as co-primary rather than treating one as
a robustness check for the other, for a substantive reason specific to
this setting: the 2020--2023 window we study spans four \emph{distinct}
legalization waves (Arizona in 2021; Rhode Island, Vermont, Montana,
New Jersey, New Mexico, and New York in 2022; Connecticut, Maryland,
and Missouri in 2023), so ``distance to the nearest legal border''
is not one clean before/after event but a continuously shifting
frontier (Figure~\ref{fig:rollout}). The continuous-distance
specification uses all of this variation; the binary switcher
specification collapses it to a single threshold crossing, which is
cruder but matches the design most directly comparable to
\citet{HaoCowan2020} and the companion arrest-substitution paper.
Reporting both, rather than picking one, is itself informative here:
where they agree, we have more confidence; where they disagree (as
they do for our main overdose result), that disagreement is the
finding.

% ============================================================
% 2. DATA AND DESIGN
% ============================================================
\section{Data and Design}

\subsection{Treatment: Two Co-Primary Specifications at a 50-Mile Radius}

We use a panel of 2,347 counties in the 32 U.S. states that had not
legalized recreational marijuana through 2023, with each county's
distance to the nearest legal-market county computed via Haversine
distance between county centroids (Census Gazetteer coordinates),
recomputed independently for this paper and validated against the
companion project's original 2020/2023 endpoint distances (mean
absolute discrepancy under 0.005 miles). We use a 50-mile
switcher/control radius throughout, tighter than the 100 miles used in
the companion arrest-substitution project.\footnote{We first attempted
a 20-mile radius, which leaves only 32 switcher counties nationally --
too few for reliable cluster-robust inference. Fifty miles yields 212.}

\textbf{Continuous specification.} The primary continuous treatment is
the raw change in distance to the nearest legal border,
\begin{equation}
  \Delta\text{Dist}_i \;=\; \text{Dist}_{i,2023} \;-\; \text{Dist}_{i,2020},
  \label{eq:distmi}
\end{equation}
negative for counties that became closer. We estimate
\begin{equation}
  \Delta Y_i \;=\; \alpha \;+\; \beta \cdot \frac{\Delta\text{Dist}_i}{50} \;+\; \varepsilon_i,
  \label{eq:distreg}
\end{equation}
scaled into 50-mile units, with a $\Delta\log(\text{distance})$
robustness check alongside it.

\textbf{Binary specification.} A county is a ``switcher'' if its
distance to the nearest legal border was at least 50 miles in 2020 and
fell below 50 miles by 2023; a ``control'' if it remained at least 50
miles away in both years. Counties within 50 miles in \emph{both}
years (54 counties) are excluded from the main sample, as in the
companion paper. We estimate
\begin{equation}
  \Delta Y_i \;=\; \alpha \;+\; \gamma \cdot \text{Switcher}_i \;+\; \varepsilon_i.
  \label{eq:switchreg}
\end{equation}
Standard errors are clustered by state throughout
\citep{CameronMiller2015, BertrandDufloMullainathan2004}. The full
panel is 212 switchers, 2,081 controls, and 54 always-near counties;
the outcome data sources below cover smaller subsets of it.

\subsection{The Legal-Market Rollout, 2020--2023}

\begin{figure}[htbp]
  \centering
  \includegraphics[width=\textwidth]{figures/fig1_rollout_maps.png}
  \caption{The Legal-Market Rollout and 50-Mile Treatment Frontier, by Year.}
  \label{fig:rollout}
\end{figure}

Figure~\ref{fig:rollout} maps the never-legal-state counties within 50
miles of a legal border at each snapshot year. The 2020 panel shows
the pre-2020 baseline legal states (Washington, Oregon, California,
Nevada, Colorado, Michigan, Massachusetts, Alaska); each subsequent
panel adds that year's newly-legal states in a darker shade and shows
the corresponding growth in treated (red-dot) counties along the new
borders. The four panels make visually clear what
Equations~\ref{eq:distreg}--\ref{eq:switchreg} handle statistically: a
Missouri county's treatment status, for instance, does not change at
all until 2023, while an Arizona-adjacent Utah county is treated
starting in 2021.

\subsection{Drug-Overdose Mortality: CDC VSRR Provisional County Data}

County-level, cause-specific mortality is not available through any
open, unauthenticated API: CDC WONDER, the relevant query tool, is not
reachable from our data-collection environment, and its underlying
mortality files require a manual query through its web interface. We
instead use the CDC's Vital Statistics Rapid Release (VSRR)
``Provisional County-Level Drug Overdose Death Counts,'' published on
the open data.cdc.gov Socrata platform, which reports 12-month-ending
drug-overdose death counts by county of residence. We take the count
for the 12 months ending December 2020 and December 2023, divide by
county population for a rate per 100,000, and take the first
difference as the outcome.

Two limitations are consequential enough to state plainly. First, per
NCHS confidentiality standards, any county-month cell with 1--9 deaths
is suppressed; this removes 1,505 of 3,144 counties nationally in 2020
and 1,425 in 2023, shrinking our regression sample to 928 counties (81
switchers, 847 controls) -- well under half the full treatment panel,
disproportionately excluding small, rural counties, which are exactly
the counties with the most identifying variation in a border-proximity
design. Second, this is \emph{provisional} surveillance data subject
to revision as death-certificate investigations close; it is not the
final vital-statistics record. We treat our overdose result as
suggestive for these reasons.

\subsection{Supplementary: Marijuana and Alcohol Arrests at the Same Threshold}

As a supplementary check, we rerun the companion project's SRS-style
FBI Crime Data Explorer arrest-rate outcomes (marijuana possession/sale,
codes 153+158; alcohol liquor-law and drunkenness, codes 270+280,
excluding DUI as not substance-specific) at this paper's 50-mile
threshold rather than the companion paper's original 100 miles, so
that the arrest evidence and the overdose evidence share an identical
treatment design. This is not, to our knowledge, an outcome combination
used elsewhere in the econ literature alongside a cross-border overdose
design, so we report it as a check on internal consistency rather than
as an independent contribution.

% ============================================================
% 3. RESULTS
% ============================================================
\section{Results}

\subsection{Drug-Overdose Mortality}

\begin{table}[htbp]
\centering
\caption{Effect on Drug-Overdose Mortality, Both Treatment Specifications}
\label{tab:overdose}
\begin{tabular}{@{}lccc@{}}
\toprule
 & Switcher (binary) & Distance, per 50 mi & Log-distance (robustness) \\
\midrule
Overdose rate per 100k & $-3.546^{**}$ (1.632) & $-0.554$ (0.340) & $+0.738$ (1.292) \\
 & $p=0.030$ & $p=0.103$ & $p=0.568$ \\
\bottomrule
\end{tabular}
\vspace{2pt}
\begin{minipage}{0.9\textwidth}
\footnotesize Note: Outcome is the change (2020--2023) in overdose
deaths per 100{,}000. SE clustered by state in parentheses.
$N=928$ non-suppressed counties (81 switchers, 847 controls).
$^{**}p<0.05$.
\end{minipage}
\end{table}

\begin{figure}[htbp]
  \centering
  \includegraphics[width=0.85\textwidth]{figures/fig2_overdose_coefplot.png}
  \caption{Effect on Drug-Overdose Mortality, Both Treatment Specifications (95\% CI).}
  \label{fig:coefoverdose}
\end{figure}

Switcher counties saw overdose death rates fall by 3.5 more per
100,000 than control counties over 2020--2023 ($p=0.030$;
Figure~\ref{fig:overdosetrend} shows both groups' raw trends) -- a
magnitude close to \citet{Denkyirah2026}'s state-level estimate of
roughly 3.5 fewer deaths per 100,000, despite the two papers using
unrelated identification strategies. However, the continuous-distance
specification, which uses strictly more of the underlying variation
(Section~2.2) and is our preferred specification on identification
grounds, is not significant ($p=0.103$), and the log-distance
robustness check flips sign entirely ($p=0.568$;
Figure~\ref{fig:scatteroverdose} plots the continuous regression
directly). A result present in the binary comparison but not the
continuous one is real evidence, but weaker evidence than a result
that holds across specifications -- combined with the sample reduction
from suppression (Section~2.3), we read this as suggestive of an
opioid-substitution channel, not confirmatory of one.

\begin{figure}[htbp]
  \centering
  \includegraphics[width=0.7\textwidth]{figures/fig3_overdose_trends.png}
  \caption{Drug-Overdose Death Rate, Switchers vs. Controls, 2020--2023.}
  \label{fig:overdosetrend}
\end{figure}

\begin{figure}[htbp]
  \centering
  \includegraphics[width=0.75\textwidth]{figures/fig4_overdose_scatter.png}
  \caption{The Continuous-Distance Regression: Distance Change vs. Overdose-Rate Change.}
  \label{fig:scatteroverdose}
\end{figure}

\subsection{Supplementary: Marijuana and Alcohol Arrests}

\begin{table}[htbp]
\centering
\caption{Supplementary: Effect on Marijuana and Alcohol Arrest Rates (50-Mile Threshold)}
\label{tab:arrests}
\begin{tabular}{@{}lccc@{}}
\toprule
 & Switcher (binary) & Distance, per 50 mi & Log-distance (robustness) \\
\midrule
Marijuana arrests & $-2.242$ (1.587) & $-0.108$ (0.382) & $+0.874$ (1.508) \\
 & $p=0.158$ & $p=0.778$ & $p=0.562$ \\
\addlinespace
Alcohol arrests & $+5.431$ (5.656) & $-3.174^{**}$ (1.413) & $-8.036^{**}$ (2.704) \\
 & $p=0.337$ & $p=0.025$ & $p=0.003$ \\
\bottomrule
\end{tabular}
\vspace{2pt}
\begin{minipage}{0.95\textwidth}
\footnotesize Note: SE clustered by state in parentheses. $N=1{,}894$
counties (157 switchers, 1,737 controls). $^{**}p<0.05$.
\end{minipage}
\end{table}

\begin{figure}[htbp]
  \centering
  \includegraphics[width=\textwidth]{figures/fig5_arrest_supp_coefplot.png}
  \caption{Supplementary: Marijuana and Alcohol Arrest Rates at the Same 50-Mile Threshold (95\% CI).}
  \label{fig:coefarrests}
\end{figure}

Marijuana arrests show no significant relationship with either
treatment. Alcohol arrests, by contrast, \emph{rise} as counties become
closer to a legal border -- significant in \emph{both} continuous
specifications (distance, $p=0.025$; log-distance, $p=0.003$), though
not in the binary switcher specification ($p=0.337$). This directly
reproduces the companion paper's headline complementarity signal (a
positive-on-closeness alcohol-arrest effect, the opposite sign from a
substitution story) at a tighter radius and a different sample, and
does so more robustly across specifications than our own overdose
result manages. We read this as evidence that the underlying treatment
design is capturing something real at the county level -- it is not
that this design generically fails to find effects -- which raises our
confidence in the overdose result's binary-specification finding even
though the continuous specification does not corroborate it.

% ============================================================
% 4. DISCUSSION
% ============================================================
\section{Discussion}

Two results in this paper are more robust than one. The alcohol-arrest
supplementary result holds across both continuous specifications
(distance and log-distance) at a substantially tighter radius than the
companion project used, which is a meaningful robustness confirmation
of that earlier finding rather than a coincidence of specification
choice. The overdose-mortality result, our primary interest, does not
clear the same bar: it is sizable and significant in the binary
switcher comparison, coincidentally close in magnitude to
\citet{Denkyirah2026}'s state-level estimate, but null in the
continuous-distance specification we otherwise prefer. Given that
\citet{Shover2019} document exactly this kind of specification- and
vintage-sensitivity in the broader medical-marijuana-opioid literature,
we do not think our overdose finding should be read as resolving the
substitution question, even directionally. It is, at minimum, the
first evidence on this specific question -- a recreational-legalization,
cross-border-distance design applied to overdose mortality -- and a
motivation for revisiting it once final (non-provisional) NCHS
mortality data for these years are released.

% ============================================================
% 5. LIMITATIONS
% ============================================================
\section{Limitations}

\begin{enumerate}[label=(\arabic*), leftmargin=*]
\item \textbf{Overdose mortality data are provisional, not final.}
CDC's VSRR county file is explicitly a rapid-release product subject
to revision; a final-data replication is not yet possible for 2023.
\item \textbf{Small-cell suppression removes more than half the panel
from the overdose regression.} Only 928 of 2,347 counties have
non-suppressed data in both years, and excluded counties are
disproportionately rural -- exactly where a border-proximity design has
the most identifying variation.
\item \textbf{The 50-mile radius is a data-driven compromise, not a
theoretically ideal catchment distance.} We could not use an initially
preferred 20-mile radius because it leaves only 32 switcher counties
nationally.
\item \textbf{Two-period design.} The switcher/control comparison uses
only 2020 and 2023, precluding a formal pre-trend test, though
Figure~\ref{fig:rollout} at least makes the staggered timing of
treatment onset within that window visually explicit.
\item \textbf{No mechanism decomposition.} We cannot distinguish
substitution in pain management, recreational substitution, reduced
enforcement freeing resources, or compositional change in who lives in
switcher counties as the specific channel behind either the overdose
or the arrest result.
\end{enumerate}

% ============================================================
% 6. CONCLUSION
% ============================================================
\section{Conclusion}

Using a cross-border continuous-distance design -- to our knowledge not
previously applied to overdose mortality -- alongside a binary
switcher comparison at a 50-mile radius, we find that counties that
became close to a legal recreational-marijuana market saw a
significant decline in drug-overdose deaths relative to counties that
did not, but only in the binary specification; the continuous-distance
result is not significant. A supplementary check using marijuana and
alcohol arrest rates at the identical threshold reproduces a companion
paper's alcohol-complementarity finding robustly across both
continuous specifications, which we take as evidence the design itself
is sound even where the overdose result is not fully conclusive. We
recommend the overdose finding be treated as motivation for a
follow-up using final mortality data, not as settled evidence of an
opioid-substitution effect.

% ============================================================
% REFERENCES
% ============================================================
\bibliographystyle{apalike}
\bibliography{references}

\end{document}
